\chapter{Prompts utilizados con modelos de lenguaje como agentes auxiliares}\label{anexo:claude_opus_4_5}

Este apéndice recoge, de forma reproducible, los dos \textit{prompts} empleados para apoyar la
evaluación del sistema descrita en el Capítulo~\ref{chapter:evaluacion}: el prompt de anotación
automática de entidades de referencia (Sección~\ref{section:anexo_prompt_anotacion}) y el prompt de
evaluación \englishText{end-to-end} mediante \englishText{LLM-as-judge}
(Sección~\ref{section:anexo_prompt_evaluacion}).

\section{Anotación automática de entidades de referencia}\label{section:anexo_prompt_anotacion}

Para construir el \textit{gold standard} de entidades clínicas sobre los 10 casos de evaluación
(Sección~\ref{section:conjunto_datos_evaluacion}) se empleó \englishText{Claude Opus 4.5} como agente de anotación
automática: a partir de los textos de referencia de cada caso de prueba, el modelo identifica
las entidades clínicamente relevantes, evitando así una anotación manual exhaustiva de los
10 casos. El listado~\ref{lst:prompt_anotacion_entidades} reproduce el prompt utilizado.

\begin{lstlisting}[
    caption={Prompt de identificación de entidades médicas de referencia (\englishText{Claude Opus 4.5}).},
    label={lst:prompt_anotacion_entidades},
    basicstyle=\scriptsize\ttfamily,
    breaklines=true,
    frame=single,
    numbers=none
]
Eres un anotador clinico encargado de construir un "gold standard" de entidades medicas a partir de
un texto en espanol, para evaluar un sistema de extraccion de entidades (NER) de uso clinico.

ENTRADA: un texto clinico en espanol (dictado o informe), proporcionado a continuacion.

TAREA: lee el texto cuidadosamente y extrae TODAS las entidades clinicamente relevantes que un
sistema NER ideal deberia detectar, clasificandolas en una de las siguientes categorias:

  - DIAGNOSTICO    : enfermedades, sindromes, problemas de salud (p. ej. "diabetes mellitus tipo 2",
                      "hipertension arterial").
  - PROCEDIMIENTO   : intervenciones, pruebas diagnosticas, tecnicas terapeuticas (p. ej.
                      "ecografia abdominal", "biopsia").
  - FARMACO         : principios activos o nombres comerciales de medicamentos (p. ej.
                      "Durogesic", "paracetamol").
  - DOSIS           : cantidad, unidad y/o pauta posologica asociada a un farmaco (p. ej.
                      "25 microgramos/hora", "cada 8 horas"). Si la dosis aparece pegada al
                      farmaco en el texto, anotala como una entidad DOSIS independiente, ligada al
                      farmaco mas cercano mediante el campo "farmaco_asociado".
  - ALERGIA         : sustancias o farmacos a los que el paciente presenta alergia o intolerancia
                      conocida.

REGLAS DE ANOTACION:
  1. Anota el "span" minimo y completo del termino clinico (p. ej. "hipertension arterial" como una
     sola entidad, no "hipertension" y "arterial" por separado).
  2. No anotes palabras genericas sin valor clinico por si solas (p. ej. "paciente", "tratamiento",
     "consulta") salvo que formen parte de un termino compuesto relevante.
  3. Si una misma entidad aparece varias veces en el texto, anotala cada vez que aparece, indicando
     la posicion (offset de caracter) de cada aparicion.
  4. Ante terminos ambiguos o coloquiales, infiere la entidad clinica mas probable segun el contexto
     y justifica brevemente tu decision en el campo "nota".
  5. No inventes entidades que no esten respaldadas por el texto. No omitas entidades relevantes por
     estar expresadas de forma indirecta o coloquial.

FORMATO DE SALIDA: devuelve EXCLUSIVAMENTE un JSON con la siguiente estructura, sin texto adicional:

{
  "entidades": [
    {
      "texto": "<texto exacto tal como aparece>",
      "categoria": "DIAGNOSTICO | PROCEDIMIENTO | FARMACO | DOSIS | ALERGIA",
      "inicio": <offset de caracter de inicio en el texto>,
      "fin": <offset de caracter de fin en el texto>,
      "farmaco_asociado": "<texto del farmaco relacionado, solo si categoria=DOSIS, si no null>",
      "nota": "<breve justificacion solo si hubo ambiguedad, si no cadena vacia>"
    }
  ]
}

TEXTO A ANOTAR:
<<<TEXT.TXT>>>
\end{lstlisting}

\section{Evaluación end-to-end asistida por IA (\textit{LLM-as-judge})}\label{section:anexo_prompt_evaluacion}

La evaluación de extremo a extremo descrita en la Sección~\ref{section:evaluacion_e2e} se realizó
proporcionando a un modelo de lenguaje, en calidad de juez automático, los cuatro artefactos
generados por el sistema para cada uno de los 10 casos de prueba (transcripción, entidades
enriquecidas, informe clínico y \textit{Bundle} FHIR) junto con el texto de referencia. El
listado~\ref{lst:prompt_evaluacion_e2e} reproduce el prompt utilizado para guiar dicha evaluación.

\begin{lstlisting}[
    caption={Prompt de evaluación end-to-end asistida por IA (\textit{LLM-as-judge}).},
    label={lst:prompt_evaluacion_e2e},
    basicstyle=\scriptsize\ttfamily,
    breaklines=true,
    frame=single,
    numbers=none
]
Eres un evaluador clinico-tecnico encargado de hacer una evaluacion end-to-end, automatica y reproducible
de una aplicacion de generacion de informes medicos a partir de voz. El pipeline es:

  audio -> ASR (transcripcion) -> NER/PLN (extraccion de entidades) -> enriquecimiento SNOMED CT
  -> traduccion a CIE-10 (ConceptMap) -> generacion de Informe.txt con LLM (llama3.2:8b) -> Fhir_Reporte.json (FHIR Bundle)

Los datos de prueba estan en:
  backend/asr-service/reports/Test1 .. Test10

Cada carpeta contiene 4 archivos:
  - text.txt          -> texto original "leido" (ground truth, lo que el locutor debia decir)
  - Reporte.txt        -> log de control de la app: por cada llamada/secuencia incluye tiempos de
                          procesamiento, transcripcion ASR, entidades NER (con etiqueta, origen,
                          confianza, posicion), enriquecimiento SNOMED (conceptId, termino preferido,
                          FSN, activo, estado de definicion) y traduccion a CIE-10 (sistema,
                          equivalencia, fuente del match)
  - Informe.txt        -> informe clinico final generado por el LLM a partir de las entidades/codigos
                          detectados, con un anexo de codificacion SNOMED
  - Fhir_Reporte.json  -> Bundle FHIR (Patient, Condition, MedicationStatement, etc.) construido a
                          partir de las mismas entidades/codigos

TAREA: evalua cada uno de los 10 casos (Test1..Test10) en las siguientes dimensiones, y luego agrega
los resultados en un informe final.

1. CALIDAD DE TRANSCRIPCION (ASR)
   - Compara la(s) transcripcion(es) dentro de Reporte.txt (seccion "Transcripcion" de cada secuencia,
     concatenadas en orden) contra text.txt (ground truth).
   - Calcula o estima el Word Error Rate (WER) o un porcentaje de similitud.
   - Senala errores relevantes (sustituciones, omisiones) especialmente si afectan terminos medicos o
     farmacos (p. ej. "Durogesic", dosis, unidades).

2. CALIDAD DE EXTRACCION DE ENTIDADES (NER/PLN)
   - A partir de text.txt, identifica que entidades clinicamente relevantes (problemas, farmacos,
     sintomas, dosis) deberia haber detectado un sistema ideal.
   - Compara contra las entidades realmente extraidas en Reporte.txt.
   - Calcula precision y recall aproximados (entidades correctas detectadas / total esperadas;
     entidades detectadas / entidades correctas).
   - Senala falsos positivos, falsos negativos, y entidades mal etiquetadas (p. ej. una palabra normal
     etiquetada como PROBLEM cuando no lo es).

3. CALIDAD DE NORMALIZACION TERMINOLOGICA (SNOMED CT / CIE-10)
   - Para cada entidad enriquecida, evalua si el conceptId SNOMED y el codigo CIE-10 asignados son
     medicamente correctos y coherentes con el contexto clinico (no solo con la palabra aislada).
   - Marca como ERROR los casos donde el mapeo es semanticamente incorrecto (p. ej. revisa si terminos
     ambiguos como "doloroso" o "analgesia" se mapearon a un concepto SNOMED que no corresponde al
     significado real en el contexto -- esto ya se observo como un problema potencial en Test1, verificalo).
   - Indica la tasa de entidades sin correspondencia SNOMED/CIE-10 ("Sin concepto coincidente").

4. CALIDAD DEL INFORME GENERADO POR EL LLM (Informe.txt)
   - Fidelidad: el informe contiene solo informacion presente en text.txt/Reporte.txt, o hay
     alucinaciones (datos inventados, farmacos/diagnosticos no mencionados)?
   - Completitud: recoge todos los hechos clinicos relevantes del texto original (diagnostico,
     tratamiento, cambios de medicacion, evolucion)?
   - Correccion clinica: la interpretacion medica es razonable (p. ej. no confundir un efecto adverso
     con un diagnostico nuevo)?
   - Estructura/formato: sigue una estructura clinica clara y consistente entre los 10 casos?
   - Coherencia con el anexo SNOMED: los codigos listados en el informe coinciden con los detectados
     en Reporte.txt?

5. CALIDAD Y VALIDEZ DEL FHIR (Fhir_Reporte.json)
   - Validez estructural basica: recursos bien formados (Patient, Condition, MedicationStatement),
     referencias `subject` validas, presencia de `resourceType`, `meta.profile`.
   - Consistencia: cada entidad detectada en Reporte.txt tiene su recurso FHIR correspondiente?
     los codigos SNOMED/CIE-10 en el Bundle coinciden con los de Reporte.txt e Informe.txt?
   - Senala entidades "huerfanas" (detectadas en Reporte.txt pero ausentes en el FHIR) o inventadas
     (presentes en FHIR sin respaldo en Reporte.txt).

6. RENDIMIENTO (a partir de los tiempos en Reporte.txt)
   - Resume tiempo total y por etapa (ASR, NER, SNOMED, ConceptMap) por caso y el promedio across los
     10 casos. Senala outliers.

METODOLOGIA
   - Procesa los 10 casos uno por uno con el mismo criterio (no cambies el rasero entre casos).
   - Se estricto pero justo: si una palabra similar foneticamente fue mal transcrita por el ASR pero el
     pipeline downstream la corrigio correctamente, anotalo como acierto parcial, no como fallo total.
   - Cuando haya ambiguedad clinica, indica tu razonamiento brevemente.

FORMATO DE SALIDA
   Genera un informe en Markdown con:
   1. Tabla resumen por caso (Test1..Test10) con una puntuacion 1-5 (o N/A) en cada una de las 5
      dimensiones (ASR, NER, Terminologia, Informe LLM, FHIR) + tiempo total de procesamiento.
   2. Para cada caso, una seccion breve listando los hallazgos concretos (errores, alucinaciones,
      mapeos incorrectos) -- no repitas contenido completo de los archivos, cita solo lo relevante.
   3. Una seccion de "Hallazgos transversales" con patrones que se repiten en varios casos (p. ej. un
      tipo de error sistematico del ASR, una etiqueta NER que falla recurrentemente, un tipo de
      alucinacion recurrente del LLM).
   4. Una seccion final de "Metricas agregadas" (WER medio, precision/recall medio de NER, % de
      entidades sin mapeo SNOMED/CIE-10, % de casos con alucinacion en el informe, tiempo medio total).
   5. Recomendaciones priorizadas (que etapa del pipeline conviene mejorar primero segun el impacto
      observado).

No modifiques ningun archivo de backend/asr-service/reports -- esta es una tarea de solo lectura y analisis.
\end{lstlisting}
